\documentclass[12pt]{article}

% هوامش
\usepackage{geometry}
\geometry{margin=0.8in}

% ألوان/رؤوس/تذييل/تباعد
\usepackage{xcolor}
\usepackage{fancyhdr}
\usepackage{setspace}
\usepackage{enumitem}

% دعم يونيكود ولغات (يستلزم XeLaTeX)
\usepackage{fontspec}
\usepackage{polyglossia}

% Math packages for bmatrix/pmatrix/align etc.
\usepackage{amsmath, amssymb, mathtools}

\setmainlanguage[numerals=western]{arabic} % أرقام 0-9
\setotherlanguage{english}

% الخطوط (بدّل الاسم إذا لزم)
\newfontfamily\arabicfont[Script=Arabic,Scale=1.05]{Amiri}
\newfontfamily\englishfont{Times New Roman}

% ترويسة/تذييل (نفس التموضع)
\pagestyle{fancy}
\fancyhf{}
\renewcommand{\headrulewidth}{0pt}
\cfoot{\textenglish{Page } \thepage}

% ماكرو الخيار كما كان
\newcommand{\choice}[3]{%
  \noindent\textbf{#1}.~#2\hfill#3\par\vspace{0.25em}
}

% فقرات
\setlength{\parskip}{0.4em}
\setlength{\parindent}{0pt}

\begin{document}

% العنوان العلوي كما في الشكل (يمين للأعلى: Model/Date)
\begin{flushright}
{\Large \textbf{{\textenglish{Algorithms -1-}}}}\\[0.25em]
{\normalsize \textenglish{Model: }{\textenglish{C}}}\\[0.1em]
{\normalsize \textenglish{Date: }{\textenglish{September 4, 2025}}}
\end{flushright}

\vspace{0.6cm}

% يسار: الجامعة؛ يمين: المدة (كما في الشكل)
\noindent
\begin{minipage}[t]{0.65\textwidth}
  \raggedright
  \textbf{{\textenglish{University of Aleppo}}}\\[0.3em]
\end{minipage}\hfill
\begin{minipage}[t]{0.3\textwidth}
  \raggedleft
  \textenglish{\textbf{Duration: }}{\textenglish{90 minuts}}
\end{minipage}

\vspace{0.9cm}

\section*{اختر الإجابة الصحيحة لكل سؤال}
\setstretch{1.25}

\noindent \textbf{\textenglish{1.}} {\textenglish{Please calculate }}\(\int f(x)dx\)

\noindent\choice{A}{{\textenglish{13}}}{}
\noindent\choice{B}{\underline{{\textenglish{14}}}}{}
\noindent\choice{C}{{\textenglish{15}}}{}
\noindent\choice{D}{{\textenglish{16}}}{}
\hfill {\textenglish{[2 points]}}\\[0.5em]

\noindent \textbf{\textenglish{2.}} \(x=3+6 \){\textenglish{ what is the value of x?}}

\noindent\choice{A}{{\textenglish{7}}}{}
\noindent\choice{B}{\underline{{\textenglish{9}}}}{}
\noindent\choice{C}{{\textenglish{8}}}{}
\noindent\choice{D}{{\textenglish{10}}}{}
\noindent\choice{E}{{\textenglish{0}}}{}
\hfill {\textenglish{[1 points]}}\\[0.5em]

\noindent \textbf{\textenglish{3.}} {\textenglish{Find the inverse of the following matrix: }}\(\begin{bmatrix} 1 & 3 & 4 \\ 5 & 2 & 1 \\ 6 & 8 & 1 \end{bmatrix}\)

\noindent\choice{A}{\underline{{\textenglish{There is no inverse for this matrix}}}}{}
\noindent\choice{B}{\[\begin{pmatrix} 1 & 2 & 3 \\ 4 & 5 & 6 \\ 7 & 8 & 9  \end{pmatrix}\]}{}
\noindent\choice{C}{{\textenglish{0}}}{}
\noindent\choice{D}{{\textenglish{0}}}{}
\hfill {\textenglish{[4 points]}}\\[0.5em]

\noindent \textbf{\textenglish{4.}} {\textenglish{Test}}

\noindent\choice{A}{{\textenglish{13}}}{}
\noindent\choice{B}{\underline{{\textenglish{14}}}}{}
\noindent\choice{C}{{\textenglish{15}}}{}
\noindent\choice{D}{{\textenglish{16}}}{}
\noindent\choice{E}{{\textenglish{17}}}{}
\hfill {\textenglish{[3 points]}}\\[0.5em]

\noindent \textbf{\textenglish{5.}} {\textenglish{Test2}}

\noindent\choice{A}{\underline{{\textenglish{1}}}}{}
\noindent\choice{B}{{\textenglish{2}}}{}
\noindent\choice{C}{{\textenglish{3}}}{}
\noindent\choice{D}{{\textenglish{4}}}{}
\noindent\choice{E}{{\textenglish{5}}}{}
\hfill {\textenglish{[3 points]}}\\[0.5em]

\noindent \textbf{\textenglish{6.}} {\textenglish{test4}}

\noindent\choice{A}{\underline{{\textenglish{1}}}}{}
\noindent\choice{B}{{\textenglish{2}}}{}
\noindent\choice{C}{{\textenglish{3}}}{}
\noindent\choice{D}{{\textenglish{4}}}{}
\noindent\choice{E}{{\textenglish{5}}}{}
\hfill {\textenglish{[1 points]}}\\[0.5em]

\noindent \textbf{\textenglish{7.}} {\textenglish{Solve for }}\(x \){\textenglish{ in the equation }}\(2x+6 = 0 \)

\noindent\choice{A}{{\textenglish{-1}}}{}
\noindent\choice{B}{{\textenglish{-2}}}{}
\noindent\choice{C}{\underline{{\textenglish{-3}}}}{}
\noindent\choice{D}{{\textenglish{3}}}{}
\hfill {\textenglish{[1 points]}}\\[0.5em]

\noindent \textbf{\textenglish{8.}} {\textarabic{احسب مشتق التابع }}\( f(x) = x^2\){\textarabic{ وأوجد نهايته عند اللانهاية }}\(\lim_{x \to \infty}f(x)\)

\noindent\choice{A}{{\textenglish{14}}}{}
\noindent\choice{B}{{\textenglish{15}}}{}
\noindent\choice{C}{\underline{{\textenglish{16}}}}{}
\noindent\choice{D}{{\textenglish{17}}}{}
\noindent\choice{E}{{\textenglish{17}}}{}
\hfill {\textenglish{[1 points]}}\\[0.5em]

\noindent \textbf{\textenglish{9.}} {\textenglish{Test3}}

\noindent\choice{A}{{\textenglish{1}}}{}
\noindent\choice{B}{{\textenglish{2}}}{}
\noindent\choice{C}{\underline{{\textenglish{3}}}}{}
\noindent\choice{D}{{\textenglish{4}}}{}
\noindent\choice{E}{{\textenglish{5}}}{}
\hfill {\textenglish{[1 points]}}\\[0.5em]

\noindent \textbf{\textenglish{10.}} {\textenglish{Hello!}}

\noindent\choice{A}{\underline{{\textenglish{1}}}}{}
\noindent\choice{B}{{\textenglish{2}}}{}
\noindent\choice{C}{{\textenglish{3}}}{}
\noindent\choice{D}{{\textenglish{4}}}{}
\hfill {\textenglish{[2 points]}}\\[0.5em]



\end{document}
